\documentclass{article}
\usepackage{amsmath}
\usepackage{xcolor}
\begin{document}

MAT 207 - ECUACIONES DIFERENCIALES

G9

ING. ANTONIO MOLINA.



\textbf{Objetivo de la Asignatura}.-  Desarrollar las habilidades suficientes para resolver problemas de modelado de la vida real en un contexto dependiente del tiempo con diferenciales, considerando el uso de Ecuaciones Diferenciales.



\textbf{Contenido}.-

1) Introduccion a las ECUACIONES DIFERENCIALES (E.D.)	1er Parcial

2) E.D.  de Primer Orden								1er Parcial

3) Aplicaciones de Modelado para E.D. de Primer Orden.		1er Parcial

4) E.D. de Orden Superior.								2do Parcial

5) Aplicaciones de Modelado para E.D. de 2do Orden.			2do Parcial

6) Series de Potencias									Ex. Final

7) Transformada de LaPlace								Ex. Final

8) Sistemas de E.D.									Ex. Final



\textbf{SISTEMA DE EVALUACION (G)}

2 EX. PARCIALES 				40%

PRACTICAS Y EX. CHICOS 		20%

Ex. FINAL					40%



\textbf{BIBLIOGRAFIA.-}

Ecuaciones Diferenciales de \textbf{Dennis Zill}

E.D. para estudiantes de Ingeniera.\textbf{ Eduardo Espinoza}

E.D. \textbf{Chungara}

Differential Equations (\textbf{Shaum})



\textbf{HORARIOS}

LUNES 		E203 	16:00 A 18:00

MIERCOLES   C102 	16:00 A 18:00    *(DIA EXAMEN -  REQUISITO \textbf{CU $\rightarrow$  1Bs})



\textbf{RED SOCIAL DE COMUNICACION}

Telegram : 



\textbf{CONSIDERACIONES EXTRA}

Se premia la participacion con:

	Resolucion de ejercicios propuestos en clases

	Responder preguntas al docente.

	Resolucion y explicacion de ejercicios en casa de forma Digital(\textbf{www.mathcha.io}) y compartirlo por el grupo.

Que son equivalentes a una firma del docente(5pts)

Puntos maximos a  ganar en clases (10 pts)



\textbf{ECAMPUS}
















\end{document}
